\documentclass[12pt]{article}

\usepackage{amsmath}
\usepackage{amssymb}
\usepackage[margin=1.25in]{geometry}
\usepackage{setspace}
\usepackage{natbib}
\usepackage{hyperref}
\usepackage{parskip}

\onehalfspacing


\title{Drawing Lines in U.S. Monetary and Fiscal history}
%\title{Reconsidering the Lines:\\
%\large Famous Reversals in Fiscal and Monetary Doctrine}

\author{Thomas J.~Sargent}
\date{\today}

\begin{document}

\maketitle

\begin{abstract}
A sequel to \citet{Sargent2011} on the trade-off between price-level
stability and Ramsey-optimal tax smoothing, this essay examines the
institutional choice between fiscal dominance and monetary dominance
through a catalogue of reversals.  Milton Friedman moved from a 1948
proposal coordinating monetary and fiscal policy to a 1960 fixed
money-growth rule; Volcker's 1979--82 attempt to implement that rule
revealed the institutional requirements it could not meet.  Clay and
Madison each opposed the First Bank on constitutional grounds, then
became architects of the Second Bank once the War of 1812 exposed the
fiscal cost of no national monetary institution.  Salmon P.\ Chase
engineered the Legal Tender Acts as Secretary of the Treasury, then as
Chief Justice wrote the opinion in \textit{Hepburn v.\ Griswold} (1870)
declaring his own Acts unconstitutional.  Three twentieth-century
parallels extend the record: Keynes moved from defender of the gold
exchange standard to author of its celebrated indictment; Eccles
championed the wartime bond-price peg and then became the leading
advocate of the Treasury-Fed Accord restoring monetary independence;
Nixon closed the gold window and declared himself a Keynesian.  Two
institutional reversals complete the picture: Congress's half-century
redrawing of the money--credit line from bond-backed National Bank Notes
to the real-bills Federal Reserve Act; and Peel's Bank Charter Act of
1844, suspended in every major crisis within a generation of its
passage.  These reversals are evidence of how contextually sensitive,
and how consequential, the institutional lines are.
\end{abstract}
%
%\tableofcontents
%\bigskip
%
%=======================================================================
\section{Introduction}
%=======================================================================

Where to draw the institutional boundaries between monetary and fiscal
authorities is an unsettled issue.
\citet{Sargent2011} argued that one source of difficulty is that 
arrangements that support price-level stability tend to foreclose the
Ramsey-optimal tax smoothing that \citet{LucasStokey1983} showed
requires using money creation as a flexible fiscal instrument, while
arrangements that enable efficient public finance tend to inject fiscal
pressures into monetary policy and thereby risk inflation.  
%There is no
%dominant option.  The choice depends on the fiscal environment,
%the quality of monetary institutions, the credibility of commitments,
%and the weight placed on different social objectives.
%
%If the choice were straightforward, one would expect thoughtful and
%principled people to make it once, correctly, and stick to it.  The
%historical record suggests otherwise. 
 Some of the ablest statesmen
and economists in American and British history changed their minds,
sometimes more than once.  This essay examines six
individual reversals in American history---Friedman, Clay, Webster,
Gallatin, Madison, Chase---and three in the twentieth century---Keynes, Eccles, Nixon---and then turns to two
institutional reversals that unfolded over decades: Congress's step by step
redrawing of the line between money and credit, from the bond-backed
National Bank Notes of the 1860s to the real-bills-backed Federal
Reserve notes of 1913; and, as a British parallel, Robert Peel's Bank
Charter Act of 1844, whose statutory line was tested in three successive
crises and bent without breaking.

Milton Friedman served as the primary example in \citet{Sargent2011}. He
moved from a 1948 proposal coordinating money creation with fiscal
deficits to a 1960 fixed $k$-percent money-growth rule that insulates money
entirely from the budget, forfeiting tax-smoothing efficiency in the
name of price-level stability.  The subsequent history of Friedman's
rule adds a further chapter: Paul Volcker's Federal Reserve
operationalized it in October 1979, only to abandon it by 1982 when
financial deregulation destabilized the very monetary aggregates on
which the rule depended.

Henry Clay and James Madison present a complementary pattern at an
earlier moment in American history.  Both opposed the First Bank of the
United States---Clay as a young Jeffersonian congressman on republican
and states'-rights grounds, Madison as leading constitutional
theorist of the founding generation.  The War of 1812, fought without a national bank capable of mobilizing
government credit, affected both of them.  Two other  conversions
illuminate the same period: Daniel Webster, a New England Federalist who
voted against the Second Bank's charter in 1816, became its most
prominent legal and political defender within three years; and Albert
Gallatin, the Jeffersonian firebrand who criticized Hamilton's Treasury
from the House floor in the 1790s, became the Bank's strongest advocate
when executive responsibility as Secretary of the Treasury convinced him of  its
fiscal indispensability.  Clay would make a national bank
a centerpiece of his ``American System''; Madison would sign the bill
creating the Second Bank in 1816, recognizing 
that long legislative practice had reframed his constitutional
concerns from 25 years earlier.

Another  individual reversal belongs to Salmon P.\ Chase.  A
hard-money Jacksonian Democrat who had 
opposed  paper currency and any  mixing of public and private
credit, Chase became Abraham Lincoln's Secretary of the Treasury in 1861
and proceeded to preside over passage of the Legal Tender Acts of 1862 and
1863, which made inconvertible paper notes (\$450 million in
``greenbacks'') legal tender for all debts.  Appointed Chief Justice
in 1864, Chase spent much of the rest of his career trying to undo his own
handiwork.  In \textit{Hepburn v.\ Griswold} (1870) he wrote the
majority opinion holding the Legal Tender Acts unconstitutional as
applied to debts contracted before their passage.  Grant's appointment  of
two new justices reversed that decision in \textit{Knox v.\ Lee} (1871).
Chase dissented.

The twentieth-century examples extend the pattern in instructive
directions.  John Maynard Keynes moved from sophisticated defender of
the gold exchange standard to an outspoken opponent, driven by
his reading of the deflationary costs that gold imposed on post-war
Britain.  His trajectory is a mirror image of Friedman's: both men
were updating the same trade-off in response to the same historical
events, but each concluded that the other's preferred failure mode was
the greater danger.  Marriner Eccles championed the wartime bond-price
peg that made the Federal Reserve a fiscal instrument of the Treasury,
then reversed course and became the leading institutional advocate for
monetary independence---at the cost of his chairmanship.  Richard
Nixon, campaigning on fiscal conservatism and sound money, closed the
gold window and embraced Keynesian stimulus in a single Sunday
announcement in August 1971.

Each reversal came with explicit accounts of why changed circumstances
warranted a new institutional line.  What the reversals collectively demonstrate is that
the stability-versus-efficiency trade-off identified in
\citet{Sargent2011} involves real
live political and constitutional choices that intelligent people have made differently in different contexts, and sometimes  differently over
the course of a single lifetime.

A coincident pattern deserves notice from the outset.  The
reversals documented here cluster around moments of systemic stress.  The
War of 1812 converted Clay and Madison simultaneously and vindicated
Gallatin's earlier administrative conclusions.  The Civil War broke
Chase and funded the bond-backed currency system.  The Great War and its
deflationary aftermath drove Keynes's evolution.  World War~II created
the circumstances that confronted  Eccles.  Nixon's 1971 reversal was in response to  fiscal consequences of Vietnam War spending.  This clustering
suggests that the ``forcing function'' is not gradual evidence-accumulation
alone, but crisis-induced revelations: costs of the existing
institutional arrangement that were tolerable in a steady fiscal
environment suddenly become undeniable to multiple actors.\footnote{''Fiscal theories of price levels'' were in the backs
of the minds of the antagonists in our story. See \citet{Cochrane2023} for a state-of-the-art presentation and critical analysis of such theories.}


%=======================================================================
\section{The Framework Recalled}
%=======================================================================

\citet{Sargent2011} builds on the \citet{SargentWallace1981}
``unpleasant monetarist arithmetic'' framework.  Let the government's
flow budget constraint be
\begin{equation}
\label{eq:flow}
b_t + \tau_t + \mu_t m_{t-1} = g_t + (1+r_{t-1})\,b_{t-1},
\end{equation}
where $b_t$ is real debt, $\tau_t$ real tax revenues, $g_t$ real spending,
$r_{t-1}$ the real interest rate on debt, and $\mu_t = (M_t/M_{t-1}-1)$
the money-growth rate applied to a real money base $m_{t-1}$.  Seigniorage
$\mu_t m_{t-1}$ enters as a source of revenue that competes with ordinary
taxation.

Two polar institutional arrangements are possible:

\textit{Monetary dominance.}  The monetary authority commits credibly to
a path for $\{M_t\}$ (or equivalently the price level), and the fiscal
authority must choose $\{\tau_t,\, b_t,\, g_t\}$ to satisfy
\eqref{eq:flow} without altering that path.  Fiscal policy is the
``residual claimant''; seigniorage is fixed exogenously by the monetary
rule.  Price-level stability is guaranteed by the commitment, but
tax-smoothing efficiency is constrained: the fiscal authority cannot use
seigniorage as a flexible instrument to smooth distortionary taxes across
states and dates.

\textit{Fiscal dominance.}  The fiscal authority sets
$\{\tau_t,\, g_t,\, b_t\}$ autonomously and the monetary authority
passively accommodates whatever seigniorage is required to clear
\eqref{eq:flow}.  Ramsey-optimal tax smoothing is now achievable:
if shocks to $g_t$ are sufficiently persistent, \citet{LucasStokey1983}
show that optimality calls for using money creation alongside other taxes.
But fiscal dominance makes the price level a fiscal consequence; a
deterioration in fiscal discipline is monetized and price-level
increases  follows.

The core result of \citet{Sargent2011} is that neither arrangement 
dominates the other.  Monetary dominance trades efficiency for stability;
fiscal dominance trades stability for efficiency.  This is the sense in which a
``line'' is not obvious: reasonable people, reasoning carefully, can end
up on different sides depending on their assessment of the relevant
fiscal and monetary environment.  The historical reversals documented
below are illustrations.


%=======================================================================
\section{Milton Friedman: From Fiscal Coordination to Monetary Rules}
%=======================================================================

Friedman's 1948 paper in the \textit{American Economic Review},
``A Monetary and Fiscal Framework for Economic Stability''
\citep{Friedman1948}, proposed an institutional architecture in which
monetary policy was explicitly subordinated to fiscal outcomes.  The
proposal rested on three pillars: (1) a requirement that commercial banks
hold 100 percent reserves against deposits, so that money creation was
entirely a government function; (2) automatic fiscal stabilizers
(progressive taxation and transfer payments) that would cause the budget
to move into surplus during booms and deficit during recessions without
discretionary action; and (3) a rule specifying that any deficit be
financed exclusively by issuing money, while any surplus be used to
retire money.  Under this system, the supply of money was mechanically
linked to the fiscal balance: the government literally printed money
to finance shortfalls and destroyed money when revenues exceeded
expenditures.

This is a case of programmatic fiscal dominance.  The monetary
instrument---the money supply---is not controlled by an independent
central bank pursuing a price-level target; it is set residually to
satisfy a fiscal accounting identity.  In the terminology of
\citet{Sargent2011}, the monetary authority has been dissolved into
the fiscal authority.  Friedman embraced this design in 1948 because he
believed that the economic cycle was predominantly driven by changes in
aggregate demand, that discretionary monetary and fiscal management were
dangerously unreliable, and that the automatic, formula-driven coupling
of money creation to fiscal outcomes would deliver macroeconomic
stabilization without giving policymakers enough discretion to do harm.
His touchstone at this stage was the \textit{coordination} of the two
policy arms, not their separation.

By the late 1950s Friedman had moved to an opposite point of view.  His Fordham University
lectures, published as \textit{A Program for Monetary Stability}
\citep{Friedman1960}, advocated for a fixed $k$-percent annual growth
rate in a well-defined monetary aggregate, entirely independent of
fiscal conditions.  The Federal Reserve should expand the money supply
at a constant rate---Friedman suggested roughly 3 to 5 percent per
year---regardless of whether the budget was in surplus or deficit.
This is a textbook case of monetary dominance: the monetary rule is the
binding commitment, and the fiscal authority must accommodate it.

Superficially, the  distance between the two positions is large.  Under
the 1948 framework, a persistent fiscal deficit would automatically
generate money creation and hence inflation; the 1960 framework
explicitly blocks this channel.  What changed?

Friedman identified several reasons \citep{Friedman1960}, Chapter 4.
First, he had become convinced by the historical evidence, which he and
\citet{FriedmanSchwartz1963} were then assembling, that autonomous
monetary shocks---changes in the money stock not driven by fiscal
conditions---were the dominant source of macroeconomic fluctuations.
The Great Depression was not a collapse of aggregate demand correctable
by fiscal-monetary coordination; it was a monetary
contraction engineered by Federal Reserve mistakes.  If monetary shocks
were the main danger, the right institutional response was to bind the
monetary authority, not to couple it to fiscal outcomes.

Second, Friedman had concluded that fiscal discipline in the US polity was not sustainable.  Governments that could finance
deficits by money creation would reliably do so.  The 1948 design assumed
that the fiscal stabilizers would be genuinely automatic and symmetric:
surpluses in good times would destroy money just as deficits in bad times
created it.  But the historical record suggested that democracies enlarged
their fiscal commitments during downturns and failed to retire them
during expansions.  A design that coupled monetary policy to fiscal
outcomes would therefore degenerate into a one-way ratchet toward
inflation.

Third, and most interesting from the perspective of
\citet{Sargent2011}, Friedman had reassessed 
the stability-efficiency trade-off.  The 1948 framework prized efficiency
(smooth automatic stabilization) and accepted the price-level risk
entailed by fiscal coupling.  The 1960 program prized stability
(a committed monetary rule) and accepted the efficiency loss from
foreclosing flexible seigniorage.  The change was not a logical
correction of an earlier error; it was a shift in the relative weight
placed on the two objectives, driven by an updated reading of what
the American political economy could actually sustain.

Friedman's reversal is a sharp illustration of the paper's main thesis:
the optimal institutional line between fiscal and monetary authority
depends on empirical and institutional features of the environment that
are themselves changing and contestable.

\subsection{Paul Volcker's experiment with  Friedman's Rule}

Friedman's $k$-percent rule was 
operationalized---at least temporarily---at the Federal Reserve.  On October 6, 1979,
Paul Volcker announced what the Fed called the ``New Operating
Procedures'': rather than targeting the federal funds rate, the Federal
Open Market Committee would henceforth target non-borrowed reserves,
in practice allowing interest rates to fluctuate freely while controlling
the growth of monetary aggregates.  The announcement was widely
interpreted, including by Friedman himself, as the closest the American
central bank had ever come to implementing the $k$-percent rule.
Volcker had not adopted Friedman's specific prescription---he was
targeting the path of money growth, not a fixed percentage---but the
animating principle of the 1979 regime was identical to the 1960
proposal: bind the monetary instrument, accept short-term
interest-rate volatility, and let the fiscal authority accommodate the
monetary path rather than the reverse.

The monetarist experiment lasted until the summer of 1982---roughly
thirty months.  What ended it was not a failure of will on Volcker's
part but an institutional problem that Friedman had identified in the
abstract yet had not fully anticipated in practice.  The financial
deregulation of the early 1980s---the Depository Institutions
Deregulation and Monetary Control Act of 1980, the proliferation of
interest-bearing checking accounts (NOW accounts), and the growth of
money-market mutual funds---fundamentally altered the demand function for
the monetary aggregates that served as the rule's target.  Velocity
became unpredictable; the historical relationship between M1 growth and
nominal income that had justified M1 as an intermediate target broke
down.  The aggregates were sending contradictory signals about the stance
of policy.

By 1982, the recession induced by the disinflation was pushing
unemployment toward 11 percent, the highest since the Great Depression.
The collapse of Penn Square Bank in July 1982 and Mexico's announcement
in August 1982 that it could not service its sovereign debt threatened
systemic financial instability.  Volcker, without formal announcement,
shifted back to interest-rate targeting.

His memoir \citep{VolckerHarper2018} reflects explicitly on the episode.
Volcker acknowledged that the monetary aggregates had become unreliable
due to financial innovation, but also maintained that the 1979--82
disinflation had accomplished its primary objective---breaking
inflationary expectations---even at the cost of an acute recession.
In the terms of \citet{Sargent2011}, the Volcker experiment confirms
the same lesson that Friedman's own reversal had taught: the optimal
institutional design depends on features of the financial environment
that change---sometimes faster than the institutions designed to manage
them.  Friedman had argued for the $k$-percent rule because he believed
the money-demand function was stable enough to make the monetary
aggregate a reliable intermediate target.  The deregulation of the 1980s
undermined that stability, validating Friedman's rule in principle while
rendering it unworkable in the specific institutional environment in
which its single attempt at implementation occurred.


%=======================================================================
\section{Henry Clay: From Republican Critic to Champion of the American
         System}
%=======================================================================

Henry Clay entered Congress as a Kentucky senator in 1806, steeped in
the Jeffersonian Republican tradition that regarded the First Bank of
the United States as an unconstitutional concentration of financial power
in private hands.  When the First Bank's twenty-year charter expired in
1811 and recharter was debated, Clay---now a newly elected member of the
House---was among those who voted against it.  The recharter bill died
narrowly: the Senate divided 17--17, and Vice President George Clinton
cast the deciding vote against; the House defeated it 65 to 64.  Clay's opposition was rooted in the Jeffersonian political economy that Andrew
Jackson would later systematize: the Bank was constitutionally suspect,
an engine of privilege, and dangerous to the republic.

Five years later, Clay was among the foremost advocates in Congress
for creating the Second Bank of the United States.  The Second Bank was
chartered in April 1816 with a twenty-year term, capitalized at \$35
million (more than three times the First Bank's capitalization), and
endowed with explicit authority to issue notes that could serve as
currency nationwide.  Clay incorporated the national bank into his
``American System''---protective tariffs, federally financed internal
improvements, and a national monetary institution---as necessary for
American development and security.

What happened between 1811 and 1816 was the War of 1812.  In the
absence of a national bank, the Madison administration found it nearly
impossible to finance the war.  State banks proliferated and issued paper
with no common standard; the government was reduced to borrowing at
punishing terms from private financiers; military operations were
hampered by the Treasury's inability to transfer funds efficiently across
the continent.  In 1814, state banks outside New England suspended specie
payments---an event that dramatized the instability of a purely
decentralized monetary system.  The government's own fiscal operations
were threatened: Treasury notes circulated at steep discounts, and the
Secretary of the Treasury reported that he could not reliably meet
the government's payroll.

Clay's reversal is intelligible as a response to this evidence.  Before
the War of 1812, the costs of a national monetary institution---the
constitutional objections, the concentration of financial privilege, the
threat to state banking interests---loomed large relative to its
benefits, which seemed speculative and ideological.  After the war, the
corresponding cost-benefit calculation ran the other way: the
demonstrated fiscal cost of the absence of a national institution was
concrete and severe, while the constitutional objections had been blunted
by the practical success of the First Bank during its twenty years of
operation.  The line between efficient national finance and the dangers
of centralized monetary power had moved, not because the underlying
principle had changed, but because the evidence on the costs of each
institutional arrangement had accumulated.

In the framework of \citet{Sargent2011}, the First Bank had been an
instrument of fiscal-monetary coordination: it held government deposits,
facilitated tax collection, managed government debt, and provided a
measure of monetary discipline through its willingness to redeem its
notes.  The two decades after its expiration amounted to an unplanned
experiment in the absence of such coordination.  Clay read the
experimental results and updated accordingly.  That is not inconsistency;
that is rationality.

The sequel to Clay's conversion is worth noting.  When Andrew Jackson
vetoed the recharter of the Second Bank in 1832 and destroyed it by
withdrawing government deposits in 1833--34, Clay emerged as the leading
opponent of Jackson's action.  The Bank War became the defining issue
of Clay's Whig Party.  The man who had voted against the First Bank in
1811 became, by the 1830s, the foremost defender of the principle of a
national monetary institution against Jacksonian hard-money democracy.
The reversal had become the fixed point from which he did not thereafter
move.

\subsection{Daniel Webster: From Opposing the Charter to Championing
             the Bank}

Daniel Webster reversed from the opposite end of the partisan spectrum.  Webster entered the House of
Representatives in 1813 as a Federalist from New Hampshire---the party
of Hamilton that had created the First Bank.  Yet when the Second Bank
bill came before the House in 1816, Webster voted against it.  His
objections were not Clay's earlier Jeffersonian ones about
constitutionality or the concentration of privilege; Webster opposed the
\emph{specific terms} of the bill, arguing that it failed to require
prompt resumption of specie payments and that the Bank's capital
structure and relationship to the government were poorly designed.  He
wanted a sounder national bank, not no national bank at all.

Within three years Webster had become the Bank's foremost legal
champion.  He argued the Bank's case before the Supreme Court in
\textit{McCulloch v.~Maryland} (1819), the decision that vindicated
the constitutionality of a national bank under the Necessary and Proper
Clause and became a cornerstone of federal implied-powers doctrine.
By the 1820s, Webster was on standing retainer as the Bank's counsel, a
relationship he maintained throughout Nicholas Biddle's presidency of
the institution.  The financial entanglement was not incidental: Webster
borrowed from the Bank and was frank about the connection, writing to
Biddle on at least one occasion to note that ``my retainer has not been
renewed or refreshed as usual'' and to inquire whether the Bank wished
the relationship to continue.  During the Bank War of the 1830s, Webster
stood beside Clay as one of the two principal Senate opponents of
Jackson's campaign to destroy the institution.

Webster's reversal is less ideological than Clay's---he moved from
objecting to a specific design to championing the institution, rather
than crossing a philosophical divide---but illustrates an additional mechanism of
conversion.  Where Clay was compelled by the evidence of fiscal
catastrophe during the War of 1812, Webster was converted by the Bank's
\emph{operational success} in its early years and by his growing
appreciation, as the leading commercial lawyer of New England, of what
a well-managed national monetary institution could accomplish for
interstate commerce and the stability of private credit.  The
institution that looked ill-designed on paper in 1816 proved
indispensable in practice by 1819.  In the framework of
\citet{Sargent2011}, Webster updated his assessment not of the
\emph{principle} of fiscal-monetary coordination but of the
\emph{institutional capacity} to execute it competently---precisely
the kind of empirical updating that makes the optimal location of the
line a moving target.


%=======================================================================
\section{James Madison: Constitutional Principle and the Doctrine of
         Liquidation}
%=======================================================================

James Madison's trajectory runs parallel to Clay's and deeper: his
opposition to the First Bank was constitutional, grounded in the theory
of his generation.

When Alexander Hamilton proposed the First Bank in January 1791,
Madison led the House opposition.  His speech of February 2, 1791 is one
of the classic statements of strict constructionism.  The Constitution
granted Congress enumerated powers.  The power to charter a corporation
was not among them.  The ``necessary and proper'' clause could not
be interpreted to extend congressional power to any means that might be
deemed convenient or useful; ``necessary'' had to mean genuinely
indispensable to the execution of an enumerated power, and chartering a
bank was not indispensable to anything the Constitution explicitly
authorized.  Madison also invoked what would later be called
the principle of original understanding: the framers'
intentions, as manifested in the ratification debates, gave no
support for so broad an implied power.

Hamilton's counter-argument won: Washington accepted Hamilton's opinion
and signed the bill.  Madison's position established the intellectual framework that Jeffersonian
Republicans would invoke for the next two decades, and that Jackson's
Democrats would recycle in the Bank War of the 1830s.

\subsection{Albert Gallatin: From Jeffersonian Critic to
             Executive-Branch Champion}

Albert Gallatin's trajectory illustrates the same conversion from the
vantage point of executive responsibility.  Gallatin arrived in Congress
in 1795 as a Jeffersonian Republican from western Pennsylvania, steeped
in the agrarian suspicion of Hamilton's financial architecture.  As the
leading Republican voice on fiscal affairs in the House, Gallatin
subjected Hamilton's Treasury operations---including the operations of
the Bank---to close scrutiny, viewing them as instruments of
Federalist political consolidation rather than as neutral fiscal
machinery.  His opposition was pragmatic rather than constitutional:
where Madison had argued that chartering a bank exceeded Congress's
enumerated powers, Gallatin objected that Hamilton's Bank served
Federalist patronage networks more than it served the public interest.

When Jefferson appointed Gallatin Secretary of the Treasury in 1801,
Gallatin brought these Jeffersonian convictions to the office that
Hamilton had designed.  Administrative experience changed his mind
rapidly.  The Bank proved to be not the engine of partisan privilege
that Gallatin had attacked from the House floor but an indispensable
instrument of day-to-day fiscal management: it held government deposits
safely, transferred funds efficiently across the continent, facilitated
customs collections at dispersed ports, and provided a stable medium for
servicing the national debt.  By the time the Bank's twenty-year charter
approached expiration, Gallatin had become its strongest advocate within
the administration.  In a formal report to Congress in 1809, he urged
recharter in terms that would have been unthinkable from the
Jeffersonian firebrand of the 1790s, arguing that the Bank was
``necessary'' for the responsible management of the government's
finances and warning that its dissolution would impose severe fiscal
costs.

The warning went unheeded.  The recharter bill died in 1811---by the
narrowest of margins, as noted in the preceding section---and the fiscal
catastrophe of the War of~1812 vindicated Gallatin's prediction in every
particular.  Gallatin, who served as a peace commissioner at Ghent in
1814, watched from abroad as the Treasury struggled with exactly the
problems he had foreseen: no reliable depository for government funds,
no mechanism for interstate fiscal transfers, no institutional anchor
for public credit.  The experience confirmed what his years at the
Treasury had already taught him: that the abstract Jeffersonian case
against a national bank could not survive contact with the operational
realities of governing a continental republic.  Gallatin's later
writings on banking and currency, notably his \textit{Considerations on
the Currency and Banking System of the United States} (1831), drew
explicitly on his executive experience to defend the principle of a
national monetary institution against both Jacksonian hard-money
populism and the proliferating disorder of unregulated state banking.

Gallatin's conversion was driven not by wartime emergency---which came
later---but by the accumulation of administrative evidence.
The War of~1812 confirmed what Gallatin already knew; it did not cause
his reversal.  In the language of \citet{Sargent2011}, Gallatin's
experience as Treasury Secretary updated his estimate of the
\emph{efficiency gains} from fiscal-monetary coordination: the
institution that had appeared, from the opposition benches, to be a
source of corruption and privilege turned out, from the Treasury desk,
to be a source of fiscal efficiency that no combination of state banks
and ad hoc arrangements could replicate.

\subsection{Madison's Reversal}

In 1815, near the end of the war and his presidency, Madison vetoed a
bank bill---not on constitutional grounds, but because the proposed
bank lacked features that would make it an effective fiscal instrument
for the government.  The veto message explicitly declined to reopen the
constitutional question, which Madison said had been settled by
congressional and executive practice over twenty-four years.  The
following year, in 1816, Madison signed the bill creating the Second
Bank of the United States.

How did the author of the 1791 speech against bank constitutionality
come to sign the bill creating a larger successor?  Madison's answer
was sophisticated and, to his critics, convenient.
He had articulated it already in the 1815 veto message: while his own
constitutional views had not changed, the Constitution's meaning was
not fixed solely by the original understanding of individual framers.
It was also determined by the accumulated legislative, executive, and
judicial practice of the republic.  A course of action repeatedly
approved by Congress, acquiesced in by successive presidents, and
never invalidated by the courts thereby acquired constitutional
legitimacy through what Madison called the process of
``liquidation''---the settlement of constitutional meaning by
institutional practice over time.  By 1816, twenty-four years of the
First Bank's operation had, in Madison's view, liquidated the
constitutional question in favor of congressional authority.

The liquidation doctrine is consequential.  It implies
that the institutional line between permissible and impermissible
concentration of fiscal-monetary power is not fixed by original text
alone; it evolves through a common-law-like process of precedent and
practice.  Institutions that were once constitutionally contested can
become constitutionally settled, and the initial objections correspondingly
lose their force.  In the terms of \citet{Sargent2011}, the boundary
between stability-promoting and efficiency-promoting arrangements is not
only analytically ambiguous; it is constitutionally labile, subject to
reinterpretation as practice accumulates.

Madison's reversal reflects the same lesson Clay drew from 1812:
the War had revealed the fiscal consequences of the absence of a national
monetary institution.  A republic that could not finance its wars
without submitting to ruinous private financiers was ill served by
constitutional principles, however impeccably derived.

The sharpest contrast to Madison is Jefferson, who never accepted the
liquidation doctrine and insisted to the end of his life that the Bank was
unconstitutional.  Jefferson's consistency and Madison's flexibility
represent the two sides of a live debate: whether institutional practice
can revise constitutional meaning, or whether original principle must
constrain practice however inconvenient the results.


%=======================================================================
\section{Salmon P.~Chase: From Hard Money to Legal Tender to
         Constitutional Revision}
%=======================================================================

The most dramatic reversal belongs
to Salmon P.\ Chase.  His career traces a complete arc: from principled
hard-money doctrine, through wartime necessity, to a judicial crusade to
undo his own executive work.  The fullest contemporary anatomy is George
Bancroft's \textit{A Plea for the Constitution of the United States
Wounded in the House of Its Guardians} \citep{Bancroft1886}.

\subsection{A Witness of Singular Authority}

Bancroft (1800--1891) brought to monetary constitutionalism not the
moralism of a hard-money partisan but the documentary authority of a man
who had spent fifty years in the primary sources of the founding era.
He had personally known Madison and Adams; had served as Secretary of
the Navy under Polk and as Minister to Prussia and Germany under Johnson
and Grant; and was 84 when, in the wake of \textit{Juilliard v.\
Greenman} (1884), he wrote his \textit{Plea}.  The book is a
historical argument that the framers at Philadelphia had
deliberately and by large majorities stripped from the federal
government any power to make paper money legal tender for private
debts---and that a sequence of Supreme Court decisions had betrayed that
design.  His central exhibit, running through every part of the work,
is Salmon P.\ Chase.

Bancroft's constitutional argument rests on what happened in Philadelphia
on August~16, 1787.  The original draft of the Constitution had given
Congress the power ``to borrow money, and emit bills on the credit of
the United States.''  Gouverneur Morris moved to strike the words ``and
emit bills.''  The motion passed 9--2.  Madison's convention notes
recorded that the majority ``thought it advisable not to insert this
power; thinking it probable that the omission would not disable the use
of public notes so far as they would be safe and proper, and would only
cut off the pretext for a paper currency, and particularly for making
the bills a tender either for public or private debts.''  Article~I,
Section~10 then explicitly forbade the states from emitting bills of
credit or making anything but gold and silver coin legal tender.  For
Bancroft, this two-sided design---the power denied to Congress by
omission, the power denied to the states by express prohibition---was
the clearest possible evidence that the framers had closed the door to
paper legal tender at every level of government.

Chase's career is, in Bancroft's telling, the story of that door being
pried open and the constitution bleeding as a result.

\subsection{The Hard-Money Democrat}

Chase began his public career in the 1830s embedded in the
Jacksonian hard-money tradition: paper money not convertible to specie
was fraud; bank-note issuance was an illegitimate privilege; sound money
meant gold and silver coin, or notes immediately redeemable in them.
Chase was a Free Soiler and then Republican by the 1850s, but his
monetary convictions remained Jacksonian: he distrusted paper, trusted
specie, and argued that mixing governmental and banking functions
corrupted both.

As Governor of Ohio (1856--60) and then Senator, Chase never departed
from these positions.  When Lincoln appointed him Secretary of the
Treasury in March 1861, Chase brought this monetary philosophy with him.
His first instinct was to finance the Civil War through taxation and
borrowing at market terms, without resort to inconvertible paper.  He
raised tariffs, promoted the new income tax, and placed bond issues
with Jay Cooke and other financial houses.  In his initial
communications to Congress, Chase explicitly rejected proposals for a
legal-tender paper currency.

\subsection{The Legal Tender Acts and the 5-20 Bonds}

By December 1861 the situation had become critical.  Banks had
suspended specie payments; the Treasury could not meet its obligations;
military expenditures exceeded what bond markets could absorb at
sustainable rates.  Chase reversed course.
He endorsed, designed, and lobbied for the Legal Tender Act
of February~25, 1862, which authorized the issuance of \$150~million in
Treasury notes---``United States Notes,'' quickly styled
``greenbacks''---that were legal tender for all debts, public and
private, but not convertible to specie.  Subsequent acts in July~1862
and March~1863 authorized additional issues, bringing the total to
approximately \$450~million.

Alongside the greenbacks, Chase issued the so-called 5-20 bonds to
finance the war: six-percent coupon bonds redeemable after five and
payable within twenty years, whose authorizing statute (the Act of
February~25, 1862) specified that interest would be paid in coin but
was entirely silent on principal.  Chase bridged this statutory gap with
reputational commitment, marketing the bonds to the public---through Jay
Cooke's retail campaign---with the implicit understanding that
principal, too, would eventually be repaid in gold.  The commitment was
extra-statutory: it rested on Chase's personal credibility as a
hard-money man, not on any legislative guarantee.  That would matter later.

Chase's arguments for the Acts were purely fiscal and pragmatic:
the government faced a choice between legal-tender paper and either
national bankruptcy or negotiated independence.  He invoked the
Revolutionary War precedent and argued the power to issue legal-tender
notes was inherent in sovereignty.  He did not pretend consistency with
his earlier views; he argued only that emergency had foreclosed the
options he would have preferred.

For Bancroft, the Revolutionary precedent was precisely the wrong one:
continental dollars were the reductio ad absurdum of paper money and the
primary inspiration for the Philadelphia framers' decision to close the
door.  To cite them as precedent was to cite the disease as justification
for a recurrence.

In the taxonomy of \citet{Sargent2011}, the Legal Tender Acts were a
decisive move toward fiscal dominance.  The government issued an
inconvertible currency and declared it legal tender; the ``monetary
policy'' embedded in this decision---the price level would be whatever
was required to make \$450~million in greenbacks cover real government
expenditures---was entirely subordinate to fiscal necessity.  Chase had,
under pressure of circumstances, crossed from the hard-money extreme of
monetary dominance (specie standard, no fiat paper) to the
fiscal-dominance extreme (inconvertible legal-tender paper issued at the
government's discretion).

\subsection{\textit{Hepburn v.\ Griswold} (1870): Chase's First
             Reversal}

Lincoln appointed Chase Chief Justice in December~1864.  For the
remainder of his life, Chase pursued a constitutional agenda centered on
reversing the monetary consequences of his own Treasury policy.

The occasion came in \textit{Hepburn v.\ Griswold} (75 U.S.\ 603),
decided February~7, 1870.  The case involved a debt of \$11,250
contracted in 1860, before the Legal Tender Act, and tendered for
payment in 1864 with greenbacks.  The creditor, Mrs.\ Hepburn, argued
that she was entitled to gold coin equivalent, since the contract
antedated the Act that made paper legal tender.  By a 4--3 vote (two
seats were vacant), the Court held for the creditor.  Chase wrote the
majority opinion.

The \textit{Hepburn} opinion rested on two grounds.  First, Chase argued
that compelling creditors to accept depreciated paper for pre-existing
gold debts deprived them of property without due process in violation
of the Fifth Amendment.  Second, he argued that making greenbacks legal
tender for such debts was not within the Necessary and Proper Clause:
Congress could permissibly issue paper currency to finance the war, but
the further step of compelling private creditors to accept it for
pre-existing obligations was not ``necessary'' to any express power in
the required sense.  Crucially, Chase appealed not to express textual
prohibition but to the ``spirit'' of the Constitution---a methodological
choice that would shortly be turned against him.

The opinion conspicuously did not address whether Chase had been correct
in 1862 to recommend the Acts.  But the logical entailment was
inescapable: if making greenbacks legal tender for pre-existing debts
was unconstitutional, then Chase had committed a constitutional error as
Secretary.  Justice Samuel Miller's dissent noted that the majority had struck down legislation that the
Chief Justice himself had recommended and defended as a member of the
executive branch.  Miller also criticized Chase's appeal to the
``spirit'' of the Constitution as ``abstract and intangible'' reasoning
that invited an unlimited judicial veto over congressional action.

\subsection{\textit{Knox v.\ Lee} (1871): The Court-Packing Reversal}

The reversal of \textit{Hepburn} came with stunning speed and in
circumstances whose constitutional impropriety has been debated ever
since.  On the very day \textit{Hepburn} was announced---February~7,
1870---President Grant nominated two new justices to fill existing
vacancies: William Strong of Pennsylvania and Joseph Bradley of New
Jersey.  Chase believed, and said, that Grant had deliberately timed
the nominations to reconstitute the Court and obtain a reversal.  When
the Court reached its new complement of nine, the Attorney General moved
for reargument of the legal tender question.  \textit{Knox v.\ Lee} and
\textit{Parker v.\ Davis} (79 U.S.\ 457) were decided together on
May~1, 1871, by a vote of 5--4, overruling \textit{Hepburn}.

Justice Strong's majority opinion relied on the implied-powers doctrine
of \textit{McCulloch v.\ Maryland}: if Congress has the war power and
the power to borrow money, and if a legal-tender paper currency is an
appropriate means of executing those powers, the Acts are constitutional
under the Necessary and Proper Clause.  Strong's opinion included a passage chiding Chase for having ``forced'' an earlier decision when
the Court was divided and on the verge of receiving new appointees---
the majority lecturing the former majority about the indecorum of the
process by which the current majority had come to exist.

Chase dissented, joined by Justices Clifford and Field.  His dissent
insisted that the majority had abandoned the constitutional limitations
on congressional power in favor of an open-ended theory of implied
powers that left no meaningful constraint on federal action.  Field's
dissent was constitutionally precise:
the doctrine that whatever is not expressly forbidden may be exercised
would ``change the whole character of our government'' from one of
limited, enumerated powers to one of unlimited powers in its relations
to the people.

Bancroft regarded the Court's composition question as corroborating his
thesis.  On a constitutional question of the first magnitude---whether
the federal government may compel private citizens to accept depreciated
paper in discharge of gold obligations---the outcome turned not on any
change in the law or the facts but on the personnel of the Court.  This was, in constitutional terms, not a
``hard constraint'' (in the sense of a textual prohibition too clear
to circumvent) but a ``soft constraint'' that yielded to political
pressure once the political composition of the appointing authority
changed.  The same logic that had made Chase's reputational commitment
on the 5-20 bond principal extra-statutory and therefore fragile had
made the constitutional prohibition on legal tender extra-textual and
therefore fragile.

\subsection{\textit{Juilliard v.\ Greenman} (1884): Legal Tender as
             Permanent Sovereignty}

The third case swept away whatever remained of the
\textit{Hepburn} principle.  In \textit{Juilliard v.\ Greenman}
(110~U.S.\ 421), decided March~3, 1884, Justice Gray's majority
opinion---joined by eight of nine justices---held that the power to
declare paper money legal tender was not limited to wartime emergencies
but was a permanent attribute of national sovereignty.  The statutory
notes at issue had originally been issued during the war, then redeemed
and reissued in peacetime under the Resumption and Refunding Act of
1878; no emergency justified them.  That did not matter.  Gray wrote:

\begin{quote}
The power to make the notes of the government a legal tender in payment
of private debts being one of the powers belonging to sovereignty in
other civilized nations, and not expressly withheld from congress by
the constitution; we are irresistibly impelled to the conclusion that
the impressing upon the treasury notes of the United States the quality
of being a legal tender in payment of private debts is an appropriate
means, conducive and plainly adapted to the execution of the undoubted
powers of congress.
\end{quote}

For Bancroft, this language was ``the language of revolution.''
The argument from what ``other civilized nations'' had done had no place
in interpreting a constitution deliberately designed to withhold powers
that every European monarchy exercised.  The argument from ``not
expressly withheld'' converted a government of enumerated powers into
one of unlimited powers: if Congress may do anything not expressly
forbidden, then the framers' enumeration of powers was a
dead letter.  Moreover, the same nine justices who joined the
\textit{Juilliard} majority had unanimously declared, just two years
earlier, that ``the government of the United States is one of delegated,
limited, and enumerated powers,'' and that ``every valid act of congress
must find in the constitution some warrant for its passage.''  Bancroft
noted that the \textit{Juilliard} opinion was therefore ``in
flagrant antagonism'' to the Court's own recent, unanimous, and
repeatedly given interpretations of the constitutional structure.

The sole dissenter was Justice Field---in Chase's \textit{Hepburn}
majority and in the \textit{Knox} dissent, now standing entirely alone.
Bancroft insisted Field was not a minority of one but the transmitted
judgment of the founding generation: beside him stood ``invincible
vouchers for the rightness of his dissent''---Wilson, Ellsworth, and
Paterson, all appointed to the Court by Washington from among the
Convention's framers.

\subsection{Constitutional Tragedy and Soft Constraints}

Bancroft read Chase's trajectory as symptomatic of a constitutional
culture that had lost its moorings, not as individual hypocrisy.  But his framing makes
the theoretical point with unusual clarity.  Chase's reputational
commitment to gold repayment of the 5-20 bond principal and his
constitutional commitment to the invalidity of legal-tender paper were
both \textit{extra-statutory}---both existed as assurances given without
explicit legislative or textual authority.  Both were undone by the same
political logic that had originally forced Chase to accept the Legal
Tender Acts as a war measure.  When the political environment changed
enough to matter, neither commitment was self-enforcing.

For the purposes of \citet{Sargent2011}, the arc of Chase's career
encapsulates the full range of the stability--efficiency trade-off.  His
early Jacksonian position was an extreme form of monetary dominance:
specie only, no government paper, the fiscal authority must adjust
entirely to monetary reality.  The Legal Tender Acts represented a
wartime capitulation to fiscal dominance: the monetary instrument
becomes whatever fiscal survival requires.  His judicial career was an
attempt to restore monetary dominance by constitutional command---to fix
in law the institutional line against which executive necessity, however
acute, could not encroach.

The \textit{Hepburn}--\textit{Knox}--\textit{Juilliard} trilogy
demonstrates that the attempt failed, and \textit{why} it failed.
The constitutional constraint Chase sought to enforce was a soft
constraint: it rested not on language expressly prohibiting the legal
tender power but on the ``spirit'' of the Constitution as Chase
interpreted it, and on the historical evidence that the Philadelphia
Convention had deliberately struck out the ``emit bills'' clause.
Bancroft thought this soft constraint should have been treated as a
hard one; the Court disagreed, and the political environment---
unwilling to foreclose a fiscal instrument it had used to win a
civil war---did not sustain the harder reading.  Institutional lines
that are not politically sustainable cannot be maintained by
constitutional command alone, and Chase's career is the clearest proof.


%=======================================================================
\section{The Money--Credit Line: From Bond-Backing to Real Bills,
         1863--1913}
\label{sec:money-credit}
%=======================================================================

The reversals catalogued above concern the boundary between monetary
and fiscal \emph{authorities}.  A related contested line---between
\emph{money} and \emph{credit}---was redrawn repeatedly by Congress
in the half-century following Chase's National Bank Acts.

\subsection{Bond-Backing and the Line at Government Debt}

The National Bank Act of 1864 drew the money--credit line at an
unusual location: United States government bonds.  National Bank
Notes---the dominant hand-to-hand currency of the post-Civil War
era---could be issued only against specified federal bonds deposited
with the Comptroller of the Currency.  That institutional officer,
whose primary function was verifying the bond portfolios, could
certify the soundness of every dollar in circulation.  The line
between what counted as ``money-eligible'' and what counted as mere
``credit'' was the difference between government bonds on one side and
every other asset---commercial loans, railroad securities, warehouse
receipts, farm mortgages---on the other.

In the framework of \citet{Sargent2011} this arrangement is a form of
\emph{monetary dominance}, but with a peculiar twist: the binding
constraint on the money supply was not a target set by an independent
central bank but the stock of outstanding federal debt.  When the
government ran the large primary surpluses of the 1870s and 1880s and
retired war bonds, the supply of eligible collateral contracted, and
the currency supply contracted with it regardless of any commercial
demand for more money.  Fiscal virtue, no less than fiscal profligacy,
could thus distort the monetary instrument; the two distorted it in
opposite directions, but neither direction was chosen for monetary
reasons.  The bond-backing requirement was a hard constraint against
inflationary over-issue, but it purchased that stability by tying the
money supply to the government's debt-management calendar rather than
to the needs of trade.

\subsection{The Inelasticity Debate as Four Decades of Disagreement
            about the Line}

Critics from the 1870s onward called the resulting currency
``inelastic.''  The agrarian South and West needed sharply higher
volumes of currency during planting and harvest seasons; the bond-backed
system offered no mechanism for such expansion.  Financial panics
punctuated the era.  Two coalitions proposed to move the
money--credit line in opposite directions.

Greenbackers and silver inflationists wanted to dissolve the
line almost entirely.  If the government simply issued more fiat
currency---greenbacks to be increased rather than retired, or silver to
be coined at a generous ratio---there was no need for any specific asset
to ``back'' the money at all.  This was, at the limit, the same fiscal
dominance that Chase had temporarily embraced during the war: the money
supply would be whatever fiscal or political convenience required.
William Jennings Bryan's ``Cross of Gold'' speech of 1896 crystallized
this position in its most direct form.

Asset-currency reformers drew the line elsewhere.  Their complaint was
not too little money but the wrong kind of backing.  They argued that
national banks should be permitted to issue notes against a wider range
of productive assets---railroad and corporate bonds, agricultural
warehouse receipts, or any ``good assets'' on a bank's balance
sheet---rather than against government bonds alone.  This was a proposal
to relocate the money--credit boundary from ``government debt only'' to
``any sufficiently sound private claim.''  By broadening the eligibility
set, they hoped, the currency supply would be free to expand with
productive activity rather than with the Treasury's debt schedule.

The Aldrich--Vreeland Act of 1908, passed in the wake of the Panic
of~1907, was a temporary halfway house.  It allowed groups of national
banks to issue emergency currency backed by ``almost any securities the
banks were holding,'' but imposed a punitive tax schedule to ensure
rapid retirement once conditions normalized.  More consequentially, it
created the National Monetary Commission, chaired by Senator Aldrich,
whose multi-volume studies of European central banking provided the
intellectual scaffolding for the Federal Reserve Act of 1913.

\subsection{The Real Bills Resolution and Its Costs}

The Federal Reserve Act resolved the inelasticity debate by drawing
the money--credit line at a new location: short-term, self-liquidating
commercial paper.  Section~13 of the Act confined Federal Reserve
discount operations to notes and bills ``arising out of actual
commercial transactions'' maturing within ninety days, and explicitly
excluded paper drawn to finance positions in ``stocks, bonds, or other
investment securities.''  The so-called \emph{real bills doctrine}---
also called the commercial loan theory of banking, with roots in Adam
Smith's \textit{Wealth of Nations} and a long subsequent
history---became the operating principle of the new institution.

The new line distinguished money from credit not by asset class but
by the \emph{purpose and duration} of the underlying transaction.  A
sixty-day bill drawn against a shipment of grain was eligible because
the transaction was real, productive, and self-liquidating: when the
grain was sold, the bill would be repaid and the money supply would
contract automatically.  A call loan financing stock market speculation
was ineligible.  In the political rhetoric of the era the line ran
between Main Street and Wall Street; in the theoretical rhetoric it
ran between ``money backed by real production'' and ``money backed
by finance.''

The real bills criterion solved the inelasticity problem: the money
supply could expand with productive activity rather than with the
federal debt stock.  But it imported a different instability, one
Henry Thornton had identified in 1802 in his
\textit{Paper Credit of Great Britain} \citep{Thornton1802}: by
accommodating every credit demand arising from commercial transactions
valued at \emph{current prices}, the criterion set up a positive
feedback from prices to money creation to prices.  In a boom, prices
rise; the nominal value of commercial transactions rises; the demand for
discount rises; and the Fed, faithfully applying the real bills
criterion, expands the money supply---adding fuel to the expansion
rather than restraining it.  The criterion is inherently procyclical.

The consequential application came in 1929--1933.  The Federal
Reserve Board member Adolph Miller launched a ``Direct Pressure'' policy:
any bank seeking discount window relief had to certify that it had never
made speculative loans.  Banks, unwilling to submit to the implied
interrogation, refrained from borrowing reserves even in the face of
withering deposit outflows.  Instead of discounting at the window, banks
failed---nearly 9{,}000 of them.  \citet{FriedmanSchwartz1963} held
the resulting one-third contraction of the money stock largely
responsible for the Great Depression.  Where the bond-backed currency
of the national banking era had been passively contractionary in the
face of fiscal surpluses, the real bills currency proved passively
contractionary in the face of banking panics.  Each line carried its
own characteristic failure mode.

\subsection{The Theoretical Afterlife and the Sargent--Wallace
            Reconsideration}

The real bills debate connects directly to the framework of
\citet{Sargent2011}.  In a 1982 article, \citet{SargentWallace1982}
showed that within overlapping-generations models the real bills
prescription---unfettered private intermediation against productive
assets---can be Pareto superior to the quantity-theory alternative,
which imposes restrictions on intermediation in order to control the
price level.  This is a version of the stability--efficiency trade-off
that \citet{Sargent2011} identifies in the fiscal--monetary domain:
the quantity-theory regime purchases price-level stability at the cost
of constrained intermediation; the real bills regime allows efficient
intermediation at the cost of weaker price-level control.
\citet{Laidler1984} objected that the Sargent--Wallace version had
abstracted away the key feature of the historical doctrine---the
restriction to specific collateral types---but the theorem demonstrated
that the money--credit distinction is not merely an institutional
accident or a historical curiosity.  It reflects a genuine analytical
trade-off whose terms depend on the structure of the economy and the
quality of the regulatory environment.

The congressional progression from bond-backed currency to real
bills---from the National Bank Acts of 1863--64 to the Federal Reserve
Act of 1913, through four decades of inelasticity complaint and
asset-currency agitation---thus illustrates in miniature the same
dynamic as the individual reversals catalogued in this paper.  The
actors who drafted the Federal Reserve Act were consciously reversing
the institutional design that Chase had supervised as Secretary of the
Treasury.  Their reversal solved one class of problems while creating
another; and the real bills doctrine embedded in the new Act would
itself require substantial revision in the Banking Acts of 1932--35,
a revision precipitated by the very failure mode that Thornton had
foreseen in 1802.  The half-century arc from National Bank Notes to
Federal Reserve notes is another turn of the same wheel: each new
institutional line is drawn in response to the failures of the previous
one, and each generates its own successor crisis.

\subsection{Peel's Bank Charter Act and Its Three Suspensions:
            A British Parallel}
\label{subsec:peel}

The dynamic of statutory lines tested and bent by crisis has a precise
British parallel in the Bank Charter Act of 1844.  Robert Peel's Act
divided the Bank of England into two departments: an Issue Department,
which could issue notes only against gold or a fixed fiduciary limit of
\pounds 14 million; and a Banking Department, which could lend freely
but had no special link to note issuance.  The Currency School theorists
who designed this structure---Samuel Jones Loyd (Lord Overstone) was
their principal theorist---intended it as a strict statutory
constraint on money creation \citep{Fetter1965}.  By tying the note
issue to gold and banning discretionary accommodation, they believed they
had drawn the line against inflationary over-issue once and for all.

The Act was suspended by Treasury letter in October~1847,
November~1857, and May~1866---in every major financial panic within a
generation of its passage.  In each case the mechanism was the same:
the Bank of England, constrained by the Act from expanding its note
issue, faced a choice between watching the banking system collapse and
requesting from the government an indemnity letter authorizing temporary
breach of the statutory limit.  The government wrote.  In~1847, the
announcement alone proved sufficient to calm the panic; the Bank's note
issue never in fact breached the legal maximum.  In~1857, following the
failure of the Western Bank of Scotland and the Glasgow City Bank, the
Bank did exceed its fiduciary limit and Parliament subsequently provided
indemnity.  In~1866, after the collapse of Overend, Gurney
and Company---then the largest bill broker in London---the Bank
initially hesitated to lend freely, exacerbating the panic before the
suspension letter arrived.

Walter Bagehot's response in \textit{Lombard Street} \citep{Bagehot1873}
is the consequential commentary on this
pattern.  Bagehot was not an opponent of the Bank Charter Act on
Currency School grounds; he accepted its basic institutional structure
as too politically entrenched to repeal.  But his prescription inverted
Peel's design.  Where Peel's Currency School architects had intended the
gold constraint to force contraction in a crisis---compelling banks to
liquidate and prices to fall, rather than accommodating speculative
excess---Bagehot argued that in a genuine panic the Bank \emph{must}
lend freely at a penalty rate against good collateral.  The lender-of-
last-resort function that \textit{Lombard Street} articulated was the
theoretical mirror image of everything the Currency School had built the
Act to prevent.

In the terms of \citet{Sargent2011}, the Bank Charter Act was an attempt
to inscribe monetary dominance into statute: the note issue would be
bound, the banking or fiscal residual would adjust.  The three
suspensions demonstrated that the attempt produced a statutory line
rather than an institutional reality.  The line was not changed; it was,
as Bagehot saw, managed by the publicly known expectation that it would
be suspended whenever truly tested.  ``The practical consequence is
that the Bank of England is obliged to keep the reserve for the whole
banking system,'' Bagehot wrote, ``and yet no one has said it should;
and no one will say it should.''  The constraint existed on paper; the
practice had already moved past it.

The parallel to the \textit{Hepburn}--\textit{Knox}--\textit{Juilliard}
trilogy is exact.  In both cases a monetary rule was established by law
or constitutional convention to prevent fiscal or speculative
accommodation.  In both cases the rule was suspended under pressure---by
executive letter in Britain, by court-packing in America.  In both cases
the exception tended to become the norm: repeated Bank Act suspensions
made suspension a publicly anticipated contingency, just as
\textit{Juilliard} made legal tender a permanent attribute of sovereignty.
And in both cases the analytical response was not a
defense of the original line but a new doctrine that institutionalized
the exception: Bagehot created the lender-of-last-resort doctrine to
rationalize the accommodation; Bancroft could only mourn that the
Philadelphia framers' judgment had been overridden.  The difference in
outcome reflects a difference in institutional creativity: Britain
produced a reformulation that preserved the spirit of stability while
accepting the necessity of crisis lending; America consumed the original
constraint without producing a comparably coherent successor.


%=======================================================================
\section{Twentieth-Century Reversals}
\label{sec:twentieth}
%=======================================================================

Three twentieth-century episodes extend the same pattern and illuminate
dimensions of the stability--efficiency trade-off that the earlier
record leaves implicit.

\subsection{John Maynard Keynes: From Gold Standard Defender to
             Its Foremost Critic}

Among economists, the sharpest reversal of the twentieth century belongs
to Keynes.  Its direction is the mirror image of Friedman's: where
Friedman moved from fiscal-monetary coordination toward a strict
monetary rule, Keynes moved from defense of a specie-based constraint
toward managed money and fiscal flexibility.

Keynes's starting point was \textit{Indian Currency and Finance}
\citep{Keynes1913}.  The book analyzed India's gold exchange standard
affirmatively: the system was a rational institutional evolution that
captured the price-stability benefits of gold convertibility without
requiring a full gold circulation.  The book was received as evidence
of a youthful attachment to specie discipline \citep{Moggridge1992}.

Within a decade the world that book described had been destroyed by the
Great War.  Britain had suspended gold convertibility in 1914; the United
States had accumulated the bulk of the world's monetary gold; wartime
inflation had restructured debt relationships in ways that made the
restoration of pre-war parities a program of deflation rather than of
price stability.  By 1923, when Keynes published \textit{A Tract on Monetary Reform}
\citep{Keynes1923}, his position had inverted.  The gold standard was
``already a barbarous relic''; the objective was to manage the price
level directly through central bank discretion rather than subordinate
domestic conditions to an external specie constraint.  The book proposed
a ``managed currency'' targeting price-level and output stability.  In
the taxonomy of \citet{Sargent2011}, this is a shift from monetary
dominance toward managed discretion---the monetary authority responding
to economic conditions rather than a fixed specie rule.

The break became concrete in ``The Economic Consequences of
Mr.~Churchill'' \citep{Keynes1931}, originally published in 1925.
Keynes attacked Churchill's decision to return Britain to gold at the
pre-war parity of \$4.86, which overvalued sterling by roughly
10~percent and required British wages and prices to fall by that margin.
The consequences followed: industrial contraction, the coal dispute that
became the General Strike of 1926, prolonged unemployment in the export
industries.  The episode confirmed that the price of monetary discipline
under gold could be measured in unemployment rates and broken industries.

Keynes's reversal mirrors Friedman's in structure while inverting its
conclusion.  Both updated the same trade-off in response to the same
cluster of events---the collapse of the pre-war monetary order, the
Great War, the interwar deflation---but each concluded that the other's
preferred failure mode was the greater danger.  Keynes, observing from
deflationary Britain, concluded that adjustment under specie discipline
was the dominant risk.  Friedman, composing the \textit{Monetary
History} with Schwartz \citep{FriedmanSchwartz1963}, concluded that
discretionary monetary contraction was the dominant risk.  Each was
correct about the failure mode he emphasized.
\citet{Eichengreen1992} shows how the gold standard constrained
monetary responses to the Depression; \citet{FriedmanSchwartz1963}
demonstrates how the Federal Reserve's discretionary failures deepened
it.  The divergent conclusions are the analytical content of the
no-dominant-option result of \citet{Sargent2011}.

\subsection{Marriner Eccles and the Treasury--Fed Accord of 1951}

The structural parallel to Chase's arc belongs to Marriner Eccles.
As Chairman of the Federal Reserve from 1934 to 1948, Eccles championed
the wartime arrangement under which the Fed pegged Treasury bond
prices---holding 90-day bill yields at $\tfrac{3}{8}$~percent and long
bond yields at $2\tfrac{1}{2}$~percent.  This was a deliberate crossing
to fiscal dominance: the monetary authority became an unlimited buyer of
government securities to keep borrowing costs low.  Eccles understood the
inflationary risk and accepted it as the price of war finance, saying so
in congressional testimony and in his memoir \citep{Eccles1951}.

As wartime controls were lifted and inflation accelerated, the peg
required the Fed to expand its balance sheet to prevent bond yields
from rising---accommodating inflationary pressure rather than resisting
it.  By 1948, Eccles's arguments for restoring Federal Reserve
independence had become a political liability.  Truman demoted him from
the chairmanship to an ordinary Board seat.  Like Chase, who used his
Chief Justiceship to attack the Acts he had designed, Eccles used his
reduced position to press for monetary independence---and was penalized
for it.  The man who had designed the fiscal-dominance arrangement was
removed from its administration precisely because he had concluded it
needed to end.

The Treasury--Fed Accord of March~4, 1951, negotiated by McCabe's
successor William Martin, formally ended the peg \citep{HetzelLeach2001}.
The Fed regained authority to set interest rates without regard to
Treasury borrowing costs.  The Accord became the founding charter of
American central bank independence \citep{Meltzer2003}.

Where Chase's \textit{Hepburn} decision was reversed by \textit{Knox}
and extinguished by \textit{Juilliard}, the Accord held.  Several
features appear decisive.  First, the inflationary consequences of the
peg were visible and politically salient: it hurt middle-class savers,
so the coalition for monetary independence extended beyond financial
elites.  Second, the Korean War generated some pressure for Treasury
accommodation but not the existential emergency of 1861--65.  Third,
the Fed had an immediately exercisable alternative---raise interest
rates without any legislative action---whereas Chase needed the Court to
strike down a statute.  The Accord held because the inflation tax imposed
by the peg had created a constituency for its removal that the Civil War
greenbacks had not.

\subsection{Richard Nixon and the End of Bretton Woods}

Nixon built his political identity as the anti-Keynesian: his 1968
campaign attacked Johnson's ``guns and butter'' policy---Vietnam War
spending combined with Great Society expansion without compensating tax
increases.  He opposed discretionary deficit spending and price controls.
His initial Council of Economic Advisers chairman, Paul McCracken, was a
monetarist-inflected fiscal conservative.

Bretton Woods---linking the dollar to gold at \$35 per ounce and other
currencies to the dollar---was the twentieth century's international
analog of the gold standard: monetary dominance at the systemic level,
with the dollar's gold peg constraining American fiscal and monetary
policy.  By~1971, persistent balance-of-payments deficits generated by
Vietnam and the Great Society had depleted the gold stock and filled
foreign central banks with dollars they were inclined to convert.  The
fiscal dominance latent in Johnson's spending had gradually undermined
the monetary constraint, as Sargent--Wallace arithmetic would predict.

On August~15, 1971, Nixon announced the New Economic Policy: he
simultaneously suspended dollar--gold convertibility, imposed a
ninety-day freeze on wages and prices, and announced a Keynesian
stimulus package designed to ensure re-election in 1972
\citep{Matusow1998}.  He told his economic advisers, ``I am now a
Keynesian in economics.''

Nixon's reversal was not compelled by military emergency.  The
balance-of-payments pressure that broke Bretton Woods was fiscal: the
accumulated legacy of ``guns and butter'' spending had inflated the
dollar monetary supply relative to gold reserves.  This is the peacetime
version of the mechanism \citet{Sargent2011} identifies: when fiscal
dominance is the revealed preference of the political system, no
external monetary anchor holds indefinitely, even without armed conflict.

Nixon announced the gold window closure as provisional.  It became
permanent.  \citet{Gowa1983} documents the political economy: the
domestic fiscal commitments that caused the balance-of-payments crisis
were not retracted, and without their retraction the gold price could
not have been defended.  The emergency of August~1971 became the
permanent monetary constitution of the world it transformed.


%=======================================================================
\section{Common Themes}
%=======================================================================

Several themes recur across these reversals.

\textit{War and systemic crisis as forcing functions---and the clustering
of reversals.}  The reversals documented here cluster at moments of acute
fiscal and financial stress.  The War of 1812 converted Clay and Madison
simultaneously, and vindicated Gallatin's earlier administrative
conclusions.  The Civil War broke Chase, constructed the bond-backed
currency system, and established the fiscal-dominance precedent that
\textit{Juilliard} would later entrench as permanent.  Britain's
participation in the Great War destroyed the international gold standard
that Keynes had analyzed affirmatively in~1913 and set in motion the
deflationary chain that drove his 1923 reversal.  World War~II created
the wartime peg that placed Eccles at the fiscal-dominance extreme;
postwar inflation drove him to the opposite pole.  Nixon's~1971 reversal
was the delayed fiscal consequence of Vietnam spending accumulated over
Eisenhower, Kennedy, and Johnson administrations.

This clustering is not coincidental.  In a steady fiscal environment,
the costs of the existing institutional arrangement are broadly distributed
across time and may be tolerable.  Under the concentrated demands of war
or systemic panic, they become immediate and undeniable.  Fiscal
dominance---using monetary instruments to finance emergency
expenditures---becomes irresistible when the alternative is defeat or
systemic collapse.  That is not an argument against institutional
constraints; it is a reminder that any constraint must be evaluated by
its behavior under the extreme stress cases that history periodically
generates.
\textit{The dynamic dimension.}
\citet{SargentWallace1981} and \citet{Sargent2011} focus on steady-state
comparisons of alternative institutional regimes.  The historical record
adds a dynamic dimension: transitions between regimes---forced by
crisis---may matter as much as the steady states themselves.  A
monetary-dominant regime that collapses under wartime pressure may leave
the fiscal authority worse off than a more flexible fiscal-dominant regime
that never pretended to prohibit seigniorage.  The value of an
institutional commitment depends not only on its steady-state properties
but on its resilience when tested.

\textit{The reflexivity of experience.}  Each reversal appealed to
experience that the original position could not have anticipated.  Clay
and Madison could not have known in 1791 or 1811 what the War of 1812
would reveal about the fiscal cost of the absence of a national monetary
institution.  Friedman could not have known in 1948 what his and
Schwartz's subsequent research would show about the monetary origins of
the Depression, or what postwar Keynesian political economy would reveal
about the systematic bias toward fiscal expansion.  Keynes could not
have known in 1913 what the Great War's fiscal demands would do to the
monetary system he had described.  Experience updated the effective
parameters of the trade-off: the relative magnitude of fiscal and
monetary risks, and the capacity of political institutions to manage them.

\textit{The constitutional versus the pragmatic argument.}  Chase's
career and Madison's liquidation doctrine reveal a tension between
constitutional arguments about where lines may be drawn and pragmatic
arguments about where they should be drawn.  The constitutional argument
tries to fix the line against the pressures of the moment; the pragmatic
argument allows the line to respond to circumstances.  Madison resolved
the tension by treating constitutional meaning as partly determined by
experience.  Chase resolved it by insisting the Constitution fixed the
line absolutely---and discovered that the political economy would not
sustain that absolutism.  Peel's Bank Charter Act illustrates the same
tension in statute: a rule rigid enough to command credibility was not
flexible enough to survive a genuine crisis, and Bagehot's
response---institutionalizing suspension as lender-of-last-resort
doctrine---was the most coherent resolution available.

\textit{The ratchet problem.}  Friedman's fixed money-growth rule was
motivated partly by his assessment that coupling monetary and fiscal
policy in a democracy reliably moves one way: toward more money
creation, more inflation, more fiscal expansion.  The record supports
this.  Clay and Madison, having supported the Second Bank, did not
seek to use it for fiscal dominance; but Jackson's destruction of the
Bank was followed by the pet-bank era, the suspension of 1837, and
eventually the Legal Tender Acts.  Nixon announced the gold window
closure as temporary; it became permanent.  Chase's emergency greenbacks
were validated by \textit{Juilliard} as a permanent attribute of
sovereignty.  The ratchet runs toward fiscal dominance when monetary
institutions are weak; only durable, credible monetary commitments can
reverse it---and such commitments are, as Eccles and Volcker both
discovered, fragile under sustained fiscal pressure.

\textit{The impossibility of personal commitment.}  Chase's career shows
that individuals cannot maintain institutional lines against political
and fiscal pressure through the force of their convictions alone.
Chase stated his monetary principles publicly for decades.  Under
sufficient pressure---the fiscal demands of a war for national
survival---he abandoned them.  Having abandoned them, he could not
restore them from the bench, because the political forces that had
driven the abandonment had not changed.  Eccles faced the same
discovery.  Institutional lines must be maintained by institutions,
not individuals; and as Madison's liquidation doctrine suggests,
institutions themselves are mutable.

\textit{The asymmetry of conclusions from the same evidence.}  The same
historical events drove Keynes and Friedman to opposite conclusions.
The interwar monetary disorder drove Keynes toward managed money and
fiscal flexibility, and drove Friedman toward a binding monetary rule.
Both conclusions were correct about the failure mode each emphasized.
This is the analytical content of the no-dominant-option result of
\citet{Sargent2011}: the correct institutional conclusion from the same
historical experience depends on a judgment about which of two genuine
risks is the greater---and that judgment can differ between equally
careful observers reasoning from identical evidence.


%=======================================================================
\section{Conclusion}
%=======================================================================

The cases examined here do not resolve the trade-off between price-level
stability and Ramsey-optimal fiscal efficiency identified in
\citet{Sargent2011}.  They confirm that the trade-off is genuine, its
terms shift with circumstances, and the optimal location of the
institutional line has no permanent answer.

Two structural patterns emerge.  First, reversals cluster at moments of
systemic stress---wars, panics, depressions---when the costs of existing
arrangements become undeniable and pressure toward fiscal dominance
becomes irresistible.  Second, reversals are rarely provisional: the
emergency measure typically becomes the permanent settlement.  The Legal
Tender Acts were validated as a standing attribute of sovereignty in
\textit{Juilliard}.  The gold window, announced as a temporary closure
in 1971, was never reopened.  Each new institutional line is drawn in
response to the failures of its predecessor and generates its own
successor crisis.  The ratchet runs toward fiscal dominance, and durable
monetary commitments are---as Eccles and Volcker both learned---fragile
under sustained fiscal pressure.

The deeper lesson is not that institutions are useless but that no
arrangement is secure.  The question of where to draw the lines is
structurally resistant to final resolution: the no-dominant-option
result of \citet{Sargent2011} is not a feature of particular parameter
values but a consequence of the underlying trade-off itself.  The
reversals documented here---spanning two centuries, two countries, and
actors as carefully reasoned as Madison, Chase, Keynes, and
Friedman---are the predictable consequence of that structure.  Anyone
who believes the problem has been permanently solved has not been paying
attention.

\bigskip

%=======================================================================
% Add the following entry to where_to_draw_lines_sequel_v2.bib:
%
% @book{Cochrane2023,
%   author    = {John H. Cochrane},
%   title     = {The Fiscal Theory of the Price Level},
%   publisher = {Princeton University Press},
%   year      = {2023},
%   address   = {Princeton, NJ}
% }

\bibliographystyle{plainnat}
\bibliography{where_to_draw_lines_sequel_v2}
%=======================================================================

\end{document}
